\begin{tikzpicture}[font=\small,>=stealth]
\draw[gray!45,rounded corners=2pt] (0,0) rectangle (14.2,8.5);
\node at (7.1,7.85) {同样 $y_0=0$，先分清中心与非中心};
\node at (7.1,7) {两图同为 $x^2/4+y^2=1$，比例相同。};
\foreach \j/\cx in {1/3.5,2/10.5}{
 \begin{scope}[shift={(\cx,4.4)},scale=1.1]
 \draw[gray,->] (-2.35,0)--(2.35,0) node[right] {$x$};
 \draw[blue!75!black,thick] (0,0) ellipse(2 and 1);
 \ifnum\j=1
  \draw[orange!85!black,thick] (-2,0)--(2,0);
  \draw[purple!80!black,thick] (-1.897366596101,-.316227766017)--(1.897366596101,.316227766017);
  \draw[gray,dashed,thick] (0,-1)--(0,1);
  \fill (0,0) circle(1.8pt) node[below right=6pt] {$O=M$};
 \else
  \draw[orange!85!black,very thick] (1,-.8660254038)--(1,.8660254038);
  \fill (1,.8660254038) circle(1.6pt) node[above right=5pt] {$A$};
  \fill (1,-.8660254038) circle(1.6pt) node[below right=5pt] {$B$};
  \fill (1,0) circle(1.8pt) node[below left=5pt] {$M(1,0)$};
 \fi
 \end{scope}}
\node at (3.5,2) {$M=O$：水平、竖直、倾斜都可以};
\node at (10.5,2) {$M\ne O$：此例只能取 $x=1$};
\node at (7.1,.75) {把 $0/0$ 当成“竖直”，会丢掉左图中的其他方向。};
\end{tikzpicture}
